% !TEX root = ../thesis.tex

\chapter{Hlavné projektové artefakty systému}\label{ch:system-artefacts}

Táto kapitola sa zameriava na identifikáciu hlavných artefaktov systému, ktoré boli rozpoznané na základe analýzy požiadaviek. Artefakty predstavujú základné objekty systému, s ktorými systém pracuje, ako aj ich základné stavy. Ich definovanie je dôležité pre pochopenie štruktúry systému a jeho správania.

Na základe predchádzajúcich kapitol boli identifikované nasledovné hlavné artefakty systému: zariadenie, filtračné zariadenie, revízia, servisný zásah, notifikácia, report a používateľ.

\section{Zariadenie}

Zariadenie predstavuje základný prvok distribučného systému vody. Ide o všeobecný objekt, ktorý zahŕňa rôzne technické komponenty, ako sú čerpacie stanice, potrubia alebo filtračné jednotky.

Zariadenie obsahuje základné informácie, ako sú identifikátor, umiestnenie, technické parametre a aktuálny stav.

Základné stavy zariadenia:
\begin{itemize}
    \item aktívne,
    \item mimo prevádzky,
    \item v údržbe.
\end{itemize}

\section{Filtračné zariadenie}

Filtračné zariadenie je špecifický typ zariadenia určený na úpravu kvality vody. Je kľúčovým prvkom z hľadiska bezpečnosti a kvality pitnej vody.

Každé filtračné zariadenie má evidované informácie o type filtra, dátume inštalácie a histórii revízií.

Základné stavy filtračného zariadenia:
\begin{itemize}
    \item v prevádzke,
    \item vyžaduje revíziu,
    \item mimo prevádzky.
\end{itemize}

\section{Revízia}

Revízia predstavuje kontrolu filtračného zariadenia vykonanú v určitom čase. Obsahuje informácie o dátume vykonania, zodpovednej osobe, výsledku kontroly a prípadných zisteniach.

Revízia môže obsahovať aj prílohy vo forme protokolov alebo dokumentov.

Základné stavy revízie:
\begin{itemize}
    \item naplánovaná,
    \item vykonaná,
    \item po termíne.
\end{itemize}

\section{Servisný zásah}

Servisný zásah predstavuje činnosť vykonanú na zariadení alebo filtri s cieľom odstrániť poruchu alebo vykonať údržbu.

Zahŕňa informácie o type zásahu, dátume, vykonávateľovi a výsledku zásahu.

Základné stavy servisného zásahu:
\begin{itemize}
    \item naplánovaný,
    \item prebieha,
    \item dokončený.
\end{itemize}

\section{Notifikácia}

Notifikácia slúži na informovanie používateľov o dôležitých udalostiach v systéme, napríklad o blížiacej sa revízii alebo o poruche zariadenia.

Notifikácie môžu byť generované automaticky systémom.

Základné stavy notifikácie:
\begin{itemize}
    \item odoslaná,
    \item doručená,
    \item prečítaná.
\end{itemize}

\section{Report}

Report predstavuje výstup systému vo forme prehľadu údajov. Môže obsahovať informácie o revíziách, poruchách, zásahoch alebo stave zariadení.

Reporty sú využívané najmä na manažérsku kontrolu a audit.

\section{Používateľ}

Používateľ predstavuje osobu, ktorá pracuje so systémom. Môže ísť o rôzne role, ako napríklad administrátor, technik alebo manažér.

Každý používateľ má definované prístupové práva a oprávnenia v systéme.
